% !TEX root = ../demo.tex
% =================================================================
%  templates/blocks.tex
%  Каталог всех цветных боксов, пометки required и окружения proof.
%
%  Общий синтаксис:
%     \begin{Имя}[<опции tcolorbox>]{<заголовок>}{<метка>} ... \end{Имя}
%  Пример: \begin{Theorem}[required]{Формула Стирлинга}{stirling}
%          ... \end{Theorem}   ->  ссылка \ref{thm:stirling}
% =================================================================

\subsection{Нумерованные боксы}

\begin{Theorem}{Название теоремы}{demo-thm}
	Нумеруется внутри section: «Теорема 1.1». Бордовая рамка.
	Ссылка на этот бокс: \verb|\ref{thm:demo-thm}| $\to$ \ref{thm:demo-thm}.
\end{Theorem}

\begin{proof}
	Окружение \verb|proof| набирает «\textbf{Доказательство.}» жирным
	и ставит $\blacksquare$ в конце.
\end{proof}

\begin{proof}[Доказательство (набросок)]
	Необязательный аргумент меняет заголовок.
\end{proof}

\begin{Lemma}{}{}
	Лемма делит счётчик с теоремой. Заголовок можно не задавать.
\end{Lemma}

\begin{Proposition}{Эквивалентные выражения для нормы}{}
	Предложение, фиолетовая рамка.
	\[
	\|\mathcal{A}\| = \sup_{\|x\|_X \le 1} \|\mathcal{A}x\|_Y
	= \sup_{x \neq 0} \frac{\|\mathcal{A}x\|_Y}{\|x\|_X}.
	\]
\end{Proposition}

\begin{Statement}{Покомпонентная непрерывность}{}
	Утверждение, синяя рамка.
\end{Statement}

\subsection{Ненумерованные боксы}

\begin{Definition}{Непрерывность отображения}{}
	Определение, охристая рамка. Номера нет — определения
	идентифицируются по названию в плашке.
\end{Definition}

\begin{Example}{}{}
	Пример, зелёная рамка.
\end{Example}

\begin{Remark}{Зачем нужны координатные функции}{}
	Замечание, серая рамка.
\end{Remark}

\begin{Corollary}{Неравенство Бесселя}{}
	Следствие, розовая рамка.
	\[
	\sum_{k=1}^{\infty} c_k^2 \le \|x\|^2.
	\]
\end{Corollary}

\subsection{Пометка обязательного пункта: \texttt{[required]}}

Необязательный аргумент \verb|[required]| добавляет утолщённую красную
левую границу и красное \textbf{(!)} у заголовка. Работает с любым боксом.

\begin{Definition}[required]{Линейный оператор}{}
	Пусть $X, Y$~--- линейные пространства над $\RR$. Отображение
	$\mathcal{A}\colon X \to Y$ называется \textbf{линейным оператором}, если
	\[
	\mathcal{A}(x + \lambda y) = \mathcal{A}x + \lambda\,\mathcal{A}y
	\qquad \forall x, y \in X,\ \forall \lambda \in \RR.
	\]
\end{Definition}

\begin{Theorem}[required]{Формула Стирлинга}{}
	\[
	n!=\sqrt{2\pi n}\left(\frac{n}{e}\right)^n\left(1+O\!\left(\frac1n\right)\right),
	\qquad n\to\infty.
	\]
\end{Theorem}

\subsection{Бокс со списком внутри}

\begin{Remark}{}{}
	Внутри бокса нормально работают списки и выключные формулы.
	\begin{enumerate}
		\item Первый пункт.
		\item Второй пункт: $\displaystyle \sum_{k=1}^{N} c_k^2 \le \|x\|^2$.
		\item Третий пункт.
	\end{enumerate}
\end{Remark}
